\appendix
\section{Appendix}\label{sec:appendix}

\subsection{Proofs of Useful Lemmas}\label{subsec:useful-proofs}

Here we provide proofs of the results in \Cref{subsec:useful-lemmas} for which we haven't found an external source.

\begin{proof}[Proof of \Cref{lem:azumanoff}]
  The proof is modeled off that in \cite{Stoeckl23}, which only addresses
  one direction. It is a straightforward blend of standard proofs of the Chernoff bound and of Azuma's inequality.
%   uses a mixture of the standard approaches for the Chernoff bound and for Azuma's inequality.
%   We only prove this elementary lemma since it's hard to find a source for. The proof is modeled off that in \cite{Stoeckl23}, which only addresses one direction.

  First, the $\ge$ direction. Choose, with foresight, $z = \ln(1 + \alpha)$.
  \begin{align*}
    \Pr&\left[ \sum_{i=1}^{t} X_i \ge (1 + \alpha) \sum_{i =1}^{t} p_i \right] \\
      &= \Pr\left[ \exp( z \sum_{i=1}^{t} X_i) \ge \exp(z (1 + \alpha) \sum_{i =1}^{t} p_i) \right] \\
      &\le \frac{\EE \exp( z \sum_{i=1}^{t} X_i)}{\exp(z (1 + \alpha) \sum_{i =1}^{t} p_i)} \\
      &\le \frac{\EE[e^{z X_1} \EE[e^{z X_2} \ldots \EE[e^{z X_t} | X_1,\ldots,X_{t-1}] \ldots | X_1] ]}{\exp(z (1 + \alpha) \sum_{i =1}^{t} p_i)} \,.
  \end{align*}
  The innermost term $\EE[e^{z X_t} | X_1,\ldots,X_{t-1}]$ is, by convexity of $e^z$, $\le p_t e^z + (1-p_t) \le e^{p_t (e^z - 1)}$; after applying this upper bound, we can factor it out and bound the $X_{t-1}$ term, and so on. Thus we continue the chain of inequalities to get:
  \begin{align*}
    \le \frac{\exp\left((e^z - 1)\sum_{i=1}^{t} p_i\right)}{\exp\left(z (1 + \alpha) \sum_{i =1}^{t} p_i\right)} = \exp\left(- ((1+\alpha) \ln (1+\alpha) - \alpha) \sum_{i =1}^{t} p_i\right) \,.
  \end{align*}
  
  For the other direction, set $z = \ln(1-\alpha)$, which is $< 0$.
  This time, $\EE[e^{z X_t} | X_1,\ldots,X_{t-1}] \le p_t e^z + (1-p_t)$ because $p_t$ is a \emph{lower bound} for $\EE[X_t | X_1,\ldots,X_{t-1}]$, and $z$ is negative. That $p_t e^z + (1-p_t) \le e^{p_t (e^z - 1)}$ still holds for negative $z$, so:
  \begin{align*}
    \Pr\left[ \sum_{i=1}^{t} X_i \le (1 - \alpha) \sum_{i =1}^{t} p_i \right] &= \Pr\left[ \exp(z\sum_{i=1}^{t} X_i) \ge \exp(z (1 - \alpha) \sum_{i =1}^{t} p_i) \right] \\
      &\le \ldots \le \frac{\exp\left((e^z - 1)\sum_{i=1}^{t} p_i\right)}{\exp\left(z (1 - \alpha) \sum_{i =1}^{t} p_i\right)} \\
      &= \exp\left(- ((1-\alpha) \ln (1-\alpha) + \alpha) \sum_{i =1}^{t} p_i\right) \,. \qedhere
  \end{align*}
\end{proof}

\begin{proof}[Proof of \Cref{lem:error-reduction-by-vote}]
  For each $i \in [p]$, let $Y_i$ be the random indicator variable for the event that $X_i \ne v$. Let $\alpha = \frac{1}{2 \delta} - 1$.The probability that $v$ is not the most common element can be bounded by the probability that it is the not the majority element; by a Chernoff bound, this is:
  \begin{align*}
    \Pr[\sum_{i\in[p]} Y_i \ge \frac{1}{2} p] &= \Pr[\sum_{i\in[p]} Y_i \ge (1 + \alpha) \delta p] \le \exp\left( - ((1+\alpha) \ln(1+\alpha) - \alpha) \delta p\right) \\
      &\le \exp\left( - 0.073 ((1+\alpha) \ln(1+\alpha) ) \delta p\right) \qquad \text{since $\alpha \ge 1/6$} \\
      &\le \exp\left( - \frac{0.073 }{2\delta} \ln\frac{1}{2\delta} \delta p\right) = \left(2\delta\right)^{0.036 p} \,. \qedhere
  \end{align*}
\end{proof}


\subsection{Mechanical Proofs for \texorpdfstring{\Cref{sec:rt-lb}}{Random Tape Lower Bound}}\label{subsec:rt-lb-mech}

\begin{proof}[Proof of \Cref{lem:rt-lb-calc-multiple}]
  Let $x$ be the left hand side of Eq. \ref{eq:rt-lb-max-min-lb}.
  The left branch of the $\min(\cdot,\cdot)$ terms in Eq. \ref{eq:rt-lb-max-min-lb} is actually unnecessary. For any integer $\lambda \ge 2$, say that
  \begin{align*}
    \left(\frac{\rl^{\lambda + 1}}{n}\right)^{\frac{2}{\lambda^2 + 3 \lambda - 2}} = \max_{k \in \NN} \left(\frac{\rl^{k + 1}}{n}\right)^{\frac{2}{k^2 + 3 k - 2}} \,.
  \end{align*}
  Then in particular,
  \begin{align*}
    \left(\frac{\rl^{\lambda + 1}}{n}\right)^{\frac{2}{\lambda^2 + 3 \lambda - 2}} \ge \left(\frac{\rl^{(\lambda -1) + 1}}{n}\right)^{\frac{2}{(\lambda-1)^2 + 3 (\lambda -1)- 2}} \ge     \left(\frac{\rl^{\lambda}}{n}\right)^{\frac{2}{\lambda^2 + \lambda - 4}} \,,
  \end{align*}
  which implies
  \begin{align*}
    n^{2 \lambda + 2} = n^{(\lambda^2 + 3 \lambda - 2) - (\lambda^2 + \lambda - 4)} \ge \rl^{\lambda \cdot (\lambda^2 + 3 \lambda - 2) - (\lambda + 1) \cdot (\lambda^2 + \lambda - 4)} = \rl^{\lambda^2 + \lambda + 4} \,,
  \end{align*}
  hence we have $\rl \le n^{\frac{2 \lambda +2}{\lambda^2 + \lambda + 4}}$. 

  On the other hand, we have
  \begin{align*}
    \rl^{1/\lambda} \ge \left(\frac{\rl^{\lambda + 1}}{n}\right)^{\frac{2}{\lambda^2 + 3 \lambda - 2}} \qquad \iff \qquad n^\lambda \le \rl^{(\lambda + 1) \lambda - \frac{\lambda^2+3\lambda - 2}{2}} = \rl^{\frac{\lambda^2 - \lambda + 2}{2}} \,,
  \end{align*}
  so the left branch of the $\min(\cdot,\cdot)$ in Eq. \ref{eq:rt-lb-max-min-lb} is only smaller when $\rl \ge n^{\frac{2 \lambda}{\lambda^2 - \lambda + 2}}$. As
  \begin{align*}
    \frac{2 \lambda +2}{\lambda^2 + \lambda + 4} \le \frac{2 \lambda}{\lambda^2 - \lambda + 2} \,,
  \end{align*}
  for all $\lambda \ge 1$, it follows that the left branch of the $\min(\cdot,\cdot)$ in Eq. \ref{eq:rt-lb-max-min-lb} is only smaller than the right when the entire term is not the maximum. Thus
  \begin{align*}
    x &\ge \max_{k \in \NN} \left(\frac{\rl^{k + 1}}{n}\right)^{\frac{2}{k^2 + 3 k - 2}} \,.
  \end{align*}

  To get a looser but more easily comprehensible lower bound, we note that $\max_{k \in \NN} \log \left(\frac{\rl^{k + 1}}{n}\right)^{\frac{2}{k^2 + 3 k - 2}}$ is piecewise linear and convex in $\log \rl$. Consequently,
  we can lower bound it using the convex function $C \frac{(\log \rl)^2}{\log n}$, where $C$ is the maximum value which satisfies the inequality at all ``corner points'' of $\max_{k \in \NN} \log \left(\frac{\rl^{k + 1}}{n}\right)^{\frac{2}{k^2 + 3 k - 2}}$. These corner points occur precisely at values of $\log \rl$ where, for some $k \ge 2$, we have:
  \begin{align*}
    \left(\frac{\rl^{k + 1}}{n}\right)^{\frac{2}{k^2 + 3 k - 2}} = \left(\frac{\rl^{(k-1) + 1}}{n}\right)^{\frac{2}{(k-1)^2 + 3 (k-1) - 2}} \,.
  \end{align*}
%   which is equivalent to: 
%   \begin{align*}
%     \left(\frac{k + 1}{k^2 + 3 k - 2} - \frac{k}{(k-1)^2 + 3 (k-1) - 2}\right) \log \rl = \left(\frac{1}{k^2 + 3 k - 2} -  \frac{1}{(k-1)^2 + 3 (k-1) - 2}\right) \log n
%   \end{align*}
  Rearranging this gives:
  \begin{align*}
    \frac{\log n}{\log \rl} = \frac{(k+1) ((k-1)^2 + 3 (k-1) - 2) - (k-1+1) (k^2 + 3 k - 2)}{((k-1)^2 + 3 (k-1) - 2) - (k^2 + 3 k - 2)} = \frac{k^2 + k + 4}{2k + 2} \,.
  \end{align*}
  so the corners occur at $\rl = n^{\frac{2k + 2}{k^2 + k + 4}}$; and at such $\rl$, we have
  \begin{align*}
    \log \left(\frac{\rl^{k + 1}}{n}\right)^{\frac{2}{k^2 + 3 k - 2}} 
    &= \left(\frac{2}{k^2 + 3 k - 2} \cdot ((k+1) \frac{2 k + 2}{k^2 + k + 4} - 1)\right) \log n = \frac{2}{k^2 + k + 4} \log n \\
      &= \frac{(\log \rl)^2}{\log n} \frac{2}{k^2 + k + 4} \left(\frac{\log n}{\log \rl}\right)^2 = \frac{(\log \rl)^2}{\log n} \frac{2}{k^2 + k + 4} \left(\frac{k^2 + k + 4}{2k + 2}\right)^2 \\
      %&=  \frac{(\log \rl)^2}{\log n} \frac{2 (k^2 + k + 4)}{(2k + 2)^2} 
      &= \frac{1}{2} \frac{(\log \rl)^2}{\log n} \frac{k^2 + k + 4}{(k+1)^2} \ge \frac{15}{32} \frac{(\log \rl)^2}{\log n} \,.
  \end{align*}
  The function $\frac{k^2 + k + 4}{(k+1)^2}$ has derivative $\frac{k-7}{(k+1)^3}$ and is minimized exactly at $k = 7$, where it has value $\frac{15}{16}$. Consequently, the value $C = \frac{15}{32}$ is the best possible.
\end{proof}

\subsection{Mechanical Proofs for \texorpdfstring{\Cref{sec:rt-ub}}{Random Tape Upper Bound}}\label{subsec:rt-ub-mech}

\begin{proof}[Proof of \Cref{lem:rt-alg-params}]
  First, we handle the case where $\ceil{\log \rl} <  \floor{2 \frac{\log(n/4)}{\log 16 \rl}}$. Then $d = \ceil{\log \rl} \le \floor{2 \frac{\log(n/4)}{\log 16 \rl}} - 1$, and $\alpha = 2$. Note that $\frac{\rl}{2^{d - 1}} \ge 1$ since $2^{d - 1} \le  2^{\ceil{\rl} - 1} \le 2\rl/2 = \rl$.
  \begin{align*}
      \prod_{i=2}^{d} b_i &= \ceil{\frac{\rl}{2^{d - 1}}} 2^{d - 2} \ge \frac{\rl}{2} = \frac{\rl}{\alpha} \\
      \prod_{i=2}^{d} b_i &=  \ceil{\frac{\rl}{2^{d - 1}}} 2^{d - 2} \le 2 \frac{\rl}{2} \le \frac{4 \rl}{\alpha} \,.
  \end{align*}
  Since $b_d = \ceil{\frac{\rl}{2^{d - 1}}} = \ceil{\frac{2 \rl}{2^{\ceil{\log \rl}}}} \le 2$,
  \begin{align*}
    \prod_{i \in [d]} w_i &= (16 \rl) \prod_{i = 2}^{d} \prod_{j=i}^{d} b_j
      \le (16 \rl) \prod_{i = 2}^{d} 2^{d - i + 1} = (16 \rl) 2^{d (d-1) / 2} \\
      &\le (16 \rl) (2^{\ceil{\log \rl}})^{(d-1) / 2} \le (16 \rl) (2 \rl)^{(d-1) / 2} \\
      &\le (16 \rl) (2 \rl)^{(\floor{2 \frac{\log (n/4)}{\log 16\rl}}-2) / 2} \\
      &\le (16 \rl) (2 \rl)^{\frac{\log (n/4)}{\log 16\rl} - 1} \\
      &\le (16 \rl)^{\frac{\log (n/4)}{\log 16\rl}} = \frac{n}{4} \le n \,.
  \end{align*}

  Second, we consider the case where $d = \floor{2 \frac{\log (n/4)}{\log(16\rl)}}$. Because $n \ge 64 \rl$, $d \ge 2$, and so
  \begin{align*}
    d = \floor{2 \frac{\log (n/4)}{\log 16 \rl}} \ge \frac{2}{3} \cdot 2 \frac{\log (n/4)}{\log 16 \rl} = \frac{\log (n/4)}{\frac{3}{4} \log 16 \rl} \ge \frac{\log (n / 4)}{\log 4 \rl} \,.
  \end{align*}
  The second inequality used that $\frac{3}{4} (4 + \log \rl) \le (2 + \log \rl)$ for $\rl \ge 4$. Consequently,
  \begin{align*}
    \alpha = \left(\frac{(4 \rl)^{d}}{n/4}\right)^{\frac{2}{d(d-1)}} \ge \left(\frac{(4 \rl)^{\frac{\log (n / 4)}{\log 4 \rl}}}{n/4}\right)^{\frac{2}{d(d-1)}} =  \left(\frac{n/4}{n/4}\right)^{\frac{2}{d(d-1)}} = 1 \,.
  \end{align*}
  We now prove Eq. \ref{eq:rt-ub-param-ge-l}. Because $\prod_{i=2}^{d-1} b_i \ge \alpha^{d - 2}$,
  \begin{align*}
    \prod_{i=2}^{d} b_i = \ceil{\frac{\rl}{\alpha^{d - 1}}} \prod_{i=2}^{d-1} b_i \ge \frac{\rl}{\alpha \prod_{i=2}^{d-1} b_i} \cdot \prod_{i=2}^{d-1} b_i = \frac{\rl}{\alpha} \,.
  \end{align*}
  For Eq. \ref{eq:rt-ub-param-le-n}, we observe that
  \begin{align*}
    d \le 2 \frac{\log(n/4)}{\log(16 \rl)}  \qquad \implies \qquad 16\rl \le (n/4)^{2/d} \qquad \implies \qquad \rl \ge \alpha^{d - 1} =\frac{(4 \rl)^2}{(n/4)^{2/d}} \,,
  \end{align*}
  and thus $\rl / \alpha^{d - 1} \ge 1$, so $b_d = \ceil{\rl / \alpha^{d -1}} \le 2 \rl / \alpha^{d - 1}$. Then since  $\prod_{i=2}^{d-1} b_i \le 2 \alpha^{d - 2}$,
  \begin{align*}
    \prod_{i=2}^{d} b_i \le \frac{2 \rl}{\alpha^{d - 1}} \prod_{i=2}^{d} b_i \le \frac{4 \rl}{\alpha \prod_{i=2}^{d-1} b_i} \cdot \prod_{i=2}^{d-1} b_i \le \frac{4 \rl}{\alpha} \,.
  \end{align*}
  Finally, we prove Eq. \ref{eq:rt-ub-param-le-n}. As noted above,
  \begin{align*}
    b_d \le \frac{2 \rl}{\alpha^{d-1}} \le \frac{4 \rl}{\alpha \prod_{i=2}^{d-1} b_j} \,.
  \end{align*}
  Applying this fact to bound the left hand side of Eq. \ref{eq:rt-ub-param-le-n} gives:
  \begin{align*}
    \prod_{i \in [d]} w_i &= (16 \rl) \prod_{i = 2}^d \prod_{j=i}^d b_j = 16 \rl (b_d)^d \prod_{i = 2}^{d-1} \prod_{j=i}^{d-1} b_j \\
      &\le 16 \rl \left(\frac{4 \rl}{\alpha \prod_{i=2}^{d-1} b_j} \right)^{d-1} \prod_{i = 2}^{d-1} \prod_{j=i}^{d-1} b_j \\
      &\le \frac{4 \cdot (4 \rl)^d}{\alpha^{d-1}} \frac{1}{\prod_{i = 2}^{d-1} \prod_{j=2}^{i-1} b_j} \\
      &\le \frac{4 \cdot(4 \rl)^d}{\alpha^{d-1}} \frac{1}{\alpha^{(d-1)(d-2)/2}} \qquad \text{since $\prod_{j=2}^{d-1} b_j \ge \alpha^{d-2}$ and $b_2 \ge b_3 \ge \ldots \ge b_{d-1}$} \\
      &= \frac{4 \cdot(4 \rl)^d}{\alpha^{d(d-1)/2}} = \frac{4 \cdot(4 \rl)^{d}}{\frac{(4 \rl)^d}{n / 4}} = n \,. \qedhere
  \end{align*}
\end{proof}




\begin{proof}[Proof of \Cref{lem:rt-ub-space-usage}]
  \Cref{alg:rt} only stores two types of data: for each $i\in[d]$, the vectors $L_i \in [w_i]^{b_i}$, and the vectors $x_i \in \{0,1\}^{b_i}$. These can be stored using $b_i \log w_i$ and $b_i$, bits respectively, for a total of:
  \begin{align*}
    \sum_{i \in [d]} b_i \log(2 w_i) &\le b_1 \log(2 w_1) + \sum_{i = 2}^{d} b_i \log(2 w_i) \\
      &\le b_1 \log(32 \rl) + \sum_{i = 2}^{d} b_i \log(2 \prod_{j=i}^{d} b_i)
      \le \sum_{i=1}^{d} b_i \log(32 \rl) \,,
  \end{align*}
  since by Eq. \ref{eq:rt-ub-param-le-2l}, $\prod_{j=i}^{d} b_i \le 4 \rl$. 
  
  We now observe that $b_d \le \ceil{\alpha}$. If $d = \ceil{\log \rl}$,
  then $\alpha = b_2 = \ldots = b_{d - 1} = 2$ and $b_d = \ceil{\rl/\alpha^{d - 1}} \le 2$. On the other hand, if $d = \floor{2 \frac{\log(n/4)}{\log(16\rl)}}$, then $\alpha = \frac{(4 \rl)^{2/(d-1)}}{(n/4)^{2/(d(d-1))}}$. We have:
  \begin{align*}
    (4\rl)^{d + 1} = (4\rl)^{\floor{2 \frac{\log(n/4)}{\log(16\rl)}} + 1} \ge (4\rl)^{2 \frac{\log(n/4)}{\log(16\rl)}} = (n/4)^2 \,,
  \end{align*}
  which implies
  \begin{align*}
    \alpha = \frac{(4 \rl)^{2/(d-1)}}{(n/4)^{2/(d(d-1))}} \ge \frac{(4 \rl)^{2/(d-1)}}{(4\rl)^{(d+1)/(d(d-1))}} = (4\rl)^{(2 - \frac{d+1}{d}) \cdot \frac{1}{d - 1}} = (4\rl)^{\frac{1}{d}} \,.
  \end{align*}
  Consequently,
  \begin{align*}
    b_d = \ceil{\frac{\rl}{\alpha^{d - 1}}} = \ceil{\frac{1}{4} \frac{4 \rl}{\alpha^{d - 1}}} \le \ceil{\frac{1}{4} (4 \rl)^{1/d}} \le (4 \rl)^{1/d} \le \alpha \,.
  \end{align*}
  
  With the bound on $b_d$, and the fact that $\alpha \ge 1$ in both cases, and that $d \le \ceil{\log \rl}$ we obtain:
  \begin{align*}
    \sum_{i \in [d]} b_i &\le \min(\rl + 1, \ceil{8 \alpha} + \ceil{3 \log 1/\delta}) + (d - 1) \ceil{\alpha} \\
      &\le \min(\rl, \ceil{3 \log 1/\delta}) + (7 + d) 2 \alpha \\
      &\le \min(\rl, \ceil{3 \log 1/\delta}) + 32 \log \rl \ceil{\frac{(4 \rl)^{2/(d-1)}}{(n/4)^{2/(d(d-1))}}} \,.
  \end{align*}
  Multiplying this last quantity by $\log(32 \rl)$ gives a space bound.

  To obtain a much weaker, but somewhat more comprehensible upper bound on $\alpha$, when $d = \floor{2 \frac{\log(n/4)}{\log(16\rl)}}$, we note that:
  \begin{align*}
    \max_{\lambda \in \NN \cap [2,\infty)} \log \frac{(4 \rl)^{2/(\lambda-1)}}{(n/4)^{2/(\lambda(\lambda-1))}} &\le \log \max_{\lambda \in \RR \cap [2,\infty)} \left(\frac{2}{\lambda - 1} \log(4\rl) - \frac{2}{\lambda(\lambda-1)} \log(n/4) \right) \\
    &\le \log \left(2 \log(4\rl) \max_{\lambda \in \RR \cap [2,\infty)} \left(\frac{1}{\lambda - 1} - \frac{1}{\lambda (\lambda -1)} \frac{\log(n/4)}{\log(4\rl)} \right)\right) \,.
  \end{align*}
  Let $\gamma = \frac{\log(n/4)}{\log(4\rl)}$; this is $\ge 1$. Let $f(x) = \frac{1}{x - 1} (1 - \frac{\gamma}{x})$. We will now prove that $\max_{x \ge 2} f(x) \le \frac{1}{2\gamma}$. We note that when $x = 2$, we have:
  \begin{align*}
    f(2) = 1 - \frac{\gamma}{2} \le \frac{1}{2\gamma} \,.
  \end{align*}
  Checking the other endpoint, we have:
  \begin{align*}
    \lim_{x \rightarrow \infty} \frac{1}{x - 1} (1 - \frac{\gamma}{x}) = 0 \,.
  \end{align*}
  Since $f(x)$ is differentiable on $[2,\infty)$, if it has a maximum other than at the endpoints, then it will occur when $\frac{d}{dx} f(x) = 0$. Solving this equation, we obtain:
  \begin{align*}
    \frac{d}{dx} f(x) = -\frac{1}{\left(x-1\right)^{2}}+\frac{\gamma\left(2x-1\right)}{\left(x\left(x-1\right)\right)^{2}} = -\frac{1}{\left(x-1\right)^{2}}\left[1-\frac{\gamma\left(2x-1\right)}{x^{2}}\right] = 0 \,,
  \end{align*}
  which is true iff $x^2 = \gamma(2x - 1)$. The solutions to the quadratic equation are
  \begin{align*}
    x = \gamma-\sqrt{\gamma\left(\gamma-1\right)} \qquad \text{and} \qquad  x = \gamma+\sqrt{\gamma\left(\gamma-1\right)} \,.
  \end{align*}
  Since $\gamma \ge 1$, the $-$ branch has $x \le 1$, which is not in $[2,\infty)$. The $+$ branch is only in $[2,\infty)$ if $\gamma \ge \frac{4}{3}$. The value of $f(x)$ in this case is:
  \begin{align*}
    f(\gamma+\sqrt{\gamma\left(\gamma-1\right)}) &= \frac{1}{\gamma+\sqrt{\gamma\left(\gamma-1\right)} - 1} \left(1 - \frac{\gamma}{\gamma+\sqrt{\gamma\left(\gamma-1\right)}} \right) \\
      &= \frac{\sqrt{\gamma(\gamma-1)}}{2 \gamma - 1 + 2 \sqrt{\gamma(\gamma-1)}} \\
      &\le \frac{1}{4 \sqrt{\gamma(\gamma-1)}}  &&\hspace{-4cm}\text{(since $\sqrt{\gamma(\gamma-1)} \le 2 \gamma - 1$ for all $\gamma \ge 1$)}
      \\
      &\le \frac{1}{2 \gamma} \,. &&\hspace{-4cm}\text{(since $\gamma \le 2 \sqrt{\gamma(\gamma-1)}$ for all $\gamma \ge \frac{4}{3}$)}
  \end{align*}
  Thus, if $f(x)$ does have a maximum in $[2,\infty)$, it is $\le \frac{1}{2\gamma}$. We conclude that $f(x) \le \frac{1}{2\gamma}$ in all cases. This proves:
  \begin{align*}
    \log \alpha \le \max_{\lambda \in \NN} \log \frac{(4 \rl)^{2/(\lambda-1)}}{(n/4)^{2/(\lambda(\lambda-1))}} 
      &\le \log\left(2 \log(4 \ell) \frac{1}{2 \frac{\log(n/4)}{\log(4\rl)}} \right) \le \log \frac{  \log(4 \ell)^2}{\log(n/4)} \,. \qedhere
  \end{align*}
\end{proof}


\subsection{Mechanical Proofs for \texorpdfstring{\Cref{sec:pd-lb}}{Pseudo-Deterministic and Random Seed Lower Bounds}}\label{subsec:pd-lb-mech}


\begin{proof}[Proof of \Cref{thm:pd-lb-intro}]
  By \Cref{lem:pd-z-lb}, we have:
  \begin{align}
    z \ge \min\left(
      \frac{\rl}{36 \log\frac{2 |S|}{\rl}},
      \frac{\rl}{17280 \log\frac{2 |S|}{\rl}} \cdot \frac{\log(1/2\delta)}{\log(64 |S|) \log\frac{2 |S|}{\rl}}
    \right) \,.\label{eq:pd-lb-restate}
  \end{align}
  By \Cref{lem:pd-s-le-2zp1}, $|S| \le \min(n, 2^{z+1}) \le \min(n, 4^z)$. We will apply this
  inequality to each branch of the minimum in Eq. \ref{eq:pd-lb-restate}. First, say that
  the left part of the minimum is larger than the right. Then $z \ge \rl / (36 \log\frac{2 |S|}{\rl})$.
  Applying $|S| \le n$ and $|S| \le 4^z$, this implies:
  \begin{align*}
    z &\ge \frac{\rl}{36 \log\frac{2 n}{\rl}} \,, \qquad \text{and} \\
    z &\ge \frac{\rl}{36 \log\frac{2 \cdot 4^z}{\rl}} \ge \frac{\rl}{36 \cdot 3 z} \qquad \implies \qquad z \ge \sqrt{\frac{\rl}{108}} \,.
  \end{align*}
  Thus:
  \begin{align}
    z \ge \max\left(\frac{\rl}{36 \log\frac{2 n}{\rl}}, \sqrt{\frac{\rl}{108}} \right) \,.\label{eq:pd-lb-case-l}
  \end{align}

  Next, say that the right side of the minimum in Eq. \ref{eq:pd-lb-restate} is larger.
  Then applying $|S| \le n$ and $|S| \le 4^z$ to that side, we get:
  \begin{align*}
    z &\ge \frac{\rl \log(1/2\delta)}{17280 (\log\frac{2 n}{\rl})^2 \log(64 n)} \,, \qquad \text{and} \\
    z &\ge \frac{\rl \log(1/2\delta)}{17280 (\log\frac{2 \cdot 4^z}{\rl})^2 \log(64 \cdot 4^z)} \ge \frac{\rl \log(1/2\delta)}{17280 \cdot 3^2 \cdot 8 z^3} \qquad \implies \qquad z \ge \left(\frac{\rl \log(1/2\delta)}{1244160}\right)^{1/4} \,.
  \end{align*}
  Thus:
  \begin{align}
    z \ge \max\left(
      \frac{\rl \log(1/2\delta)}{17280 (\log\frac{2 n}{\rl})^2 \log(64 n)},
      \left(\frac{\rl \log(1/2\delta)}{1244160}\right)^{1/4}
    \right) \,.\label{eq:pd-lb-case-r}
  \end{align}

  The minimum of the lower bounds from Eqs. \ref{eq:pd-lb-case-l} and \ref{eq:pd-lb-case-r}
  holds in all cases, so:
  \begin{align*}
    z \ge \min\left(\max\left(\frac{\rl}{36 \log\frac{2 n}{\rl}}, \sqrt{\frac{\rl}{108}}\right), \max\left(\frac{\rl \log(1/2\delta)}{17280 (\log\frac{2 n}{\rl})^2 \log(64n)},\left(\frac{\rl \log(1/2\delta)}{1244160}\right)^{1/4}\right)\right) \,.
  \end{align*}
\end{proof}


\begin{proof}[Proof of \Cref{cor:rs-lb-intro}]
  The lower bound from \Cref{thm:pd-lb-intro} for $\mif(n,t)$ with error $\delta = 1/3$, showing constants, is:
  \begin{align}
    \max\left(\frac{t \log(3/2)}{17280 (\log\frac{2 n}{t})^2 \log(64n)},\left(\frac{t \log(3/2)}{1244160}\right)^{1/4}\right) \,.\label{eq:lb-rs-pd-expr}
  \end{align}
  If the space used by an algorithm, $z$, satisfies $z \ge (\rl-1)/2$, then we tautologically have a lower bound of $(\rl-1)/2$. Otherwise, we have $2z+2 \le \rl$.
  Applying \Cref{thm:ext-rs-to-pd} gives, for the left branch of the $\max$ in Eq. \ref{eq:lb-rs-pd-expr}, with $t = \floor{\frac{\rl}{2 z + 2}} \ge \frac{1}{8 z}$:
  \begin{align*}
    z &\ge \floor{\frac{\rl}{2 z + 2}} \frac{\log(3/2)}{17280 (\log\frac{2 n}{t})^2 \log(64n)} \ge \frac{\rl}{8 z} \frac{1/2}{17280 (\log (2 n))^2 \log(64n)} \,,
  \end{align*}
  which implies
  \begin{align*}
    z \ge \sqrt{\frac{\rl}{276480 (\log (2 n))^2 \log(64n)}} \ge  \sqrt{\frac{\rl}{7741440 (\log n)^3}} \,.
  \end{align*}
  For the right branch of Eq. \ref{eq:lb-rs-pd-expr}, we obtain:
  \begin{align*}
    z \ge \left(\floor{\frac{\rl}{2 z + 2}} \frac{\log(3/2)}{1244160}\right)^{1/4} \ge \left(\frac{\rl}{8 z} \frac{1}{2 \cdot 1244160}\right)^{1/4} \,,
  \end{align*}
  which implies:
  \begin{align*}
    z^{5/4} \ge \left(\frac{\rl}{19906560}\right)^{1/4} \qquad \implies \qquad z \ge \left(\frac{\rl}{19906560}\right)^{1/5} \,.
  \end{align*}
  Combining the two lower bounds, gives:
  \begin{align}
    z \ge \max\left(\sqrt{\frac{\rl}{7741440 (\log n)^3}}, \left(\frac{\rl}{19906560}\right)^{1/5}\right) = \Omega\left(\sqrt{\frac{\rl}{(\log n)^3}} + \rl^{1/5}\right) \,. \label{eq:cor-lb-base}
  \end{align}
  This lower bound is everywhere smaller than $(\rl-1)/2$, so it is compatible with the case in which $z \ge (\rl-1)/2$. 
  
  Taking the maximum of Eq. \ref{eq:cor-lb-base} and the known random oracle lower bound for $\mif(n,\rl)$ algorithms in the static setting, \Cref{lem:ext-robust-lb}, gives:
  \begin{align*}
    z &= \Omega\left(\frac{\rl^2}{n} + \sqrt{\frac{\rl}{(\log n)^3}} + \rl^{1/5}\right)  \,. \qedhere
  \end{align*}
  
\end{proof}

\endinput
